% ═══════════════════════════════════════════════════════════════════════
%  Stratégie 2 --- XAUUSD niveaux x10. Chiffres : macros générées.
% ═══════════════════════════════════════════════════════════════════════
\section{La grille x10 et sa non-stationnarité}
\label{sec:x10_grille}

\lettrine[lines=2, findent=2pt, nindent=0pt]{L}{'intuition} des niveaux ronds
est familière et elle n'est pas absurde~: les ordres à cours limité
s'accumulent sur les prix ronds, les seuils de prise de bénéfice y sont
posés, les stops s'y regroupent. Un prix qui arrive sur un chiffre rond
rencontre donc, en principe, plus de contrepartie qu'ailleurs. La stratégie
commandée par Apogée Invest fait de cette intuition son unique objet~: elle ne
regarde que la façon dont l'or aborde les prix multiples de dix dollars.

\subsection{Définition de la grille}

La grille est la suite des multiples de 10~\$. À tout instant, le prix se
situe entre deux barreaux~: celui du dessous et celui du dessus. La stratégie
ne s'intéresse qu'à celui vers lequel le prix se dirige, et la cible d'une
position est toujours le barreau suivant dans le sens du trade.

Le calcul du barreau passe par l'arithmétique entière au millième de dollar,
qui est le pas de cotation du symbole. Ce détail n'est pas cosmétique~: un
prix stocké en virgule flottante peut valoir $4409{,}999999996$ et tomber du
mauvais côté de la grille si on le divise naïvement.

\begin{figure}[htbp]
\centering
\includegraphics[width=\linewidth]{x10_fig01_grille.png}
\caption{La grille des niveaux x10 et la zone d'approche. Le prix est surveillé
dès qu'il entre dans une bande de largeur $z \cdot ATR$ devant le barreau~;
selon que le niveau cède ou tient, la position vise le barreau suivant
(\emph{breakout}) ou le barreau précédent (\emph{reversal}).}
\end{figure}

\subsection{Le pas de grille mesuré en ATR, année par année}

Voici le constat qui gouverne tout le reste de ce rapport. Dix dollars ne
valent pas la même chose selon l'année, non pas en pouvoir d'achat, mais
\textbf{en unités de volatilité de la barre de travail}.

% GÉNÉRÉ par scripts/build_x10_report_assets.py — ne pas éditer à la main.
% Toute valeur vient de results/xau_x10/*.json.
\begin{table}[htbp]
\centering
\footnotesize
\caption{Le pas de la grille mesuré en ATR M5, année par année. La « même » stratégie vise une cible de dix-sept écarts-types de barre en début d'échantillon et de moins de quatre en fin.}
\label{tab:x10_feasibility}
\begin{tabular}{l r r r r r r r}
\toprule
\textbf{Année} & \textbf{Jours} & \textbf{Prix méd.} & \textbf{ATR M5 méd.} & \textbf{10\,\$ / ATR} & \textbf{Spread / ATR} & \multicolumn{2}{c}{\textbf{Crois. / jour}} \\
& & & & & & \textit{méd.} & \textit{moy.} \\
\midrule
2019 & $260$ & $1\,406{,}35$ & $0{,}579$ & $\mathbf{17{,}3}$ & $0{,}501$ & $7$ & $8{,}7$ \\
2020 & $261$ & $1\,775{,}01$ & $1{,}439$ & $\mathbf{7{,}0}$ & $0{,}202$ & $20$ & $21{,}1$ \\
2021 & $260$ & $1\,793{,}97$ & $1{,}037$ & $\mathbf{9{,}6}$ & $0{,}280$ & $13$ & $14{,}5$ \\
2022 & $259$ & $1\,805{,}37$ & $1{,}148$ & $\mathbf{8{,}7}$ & $0{,}253$ & $15$ & $16{,}5$ \\
2023 & $259$ & $1\,945{,}65$ & $1{,}001$ & $\mathbf{10{,}0}$ & $0{,}290$ & $13$ & $14{,}3$ \\
2024 & $261$ & $2\,380{,}75$ & $1{,}477$ & $\mathbf{6{,}8}$ & $0{,}196$ & $21$ & $21{,}0$ \\
2025 & $260$ & $3\,345{,}07$ & $2{,}637$ & $\mathbf{3{,}8}$ & $0{,}110$ & $36$ & $39{,}9$ \\
\bottomrule
\end{tabular}
\end{table}


\begin{figure}[htbp]
\centering
\includegraphics[width=\linewidth]{x10_fig02_dollars_par_atr.png}
\caption{À gauche, le pas de la grille exprimé en \code{ATR} de cinq minutes~;
à droite, le nombre médian de croisements de niveau par jour. Les deux
courbes sont le même fait, vu de deux côtés.}
\end{figure}

\begin{insightbox}
\textbf{En 2019, franchir un niveau x10 demandait \XTenDollarsPerAtrFirst{}
écarts-types de barre~; en 2025, il en faut \XTenDollarsPerAtrLast.} La
« même » stratégie ne désigne pas le même objet d'une année à l'autre. L'or
traversait la grille \XTenCrossingsFirst{} fois par jour en 2019~; il la
traverse \XTenCrossingsLast{} fois par jour en 2025. Un niveau psychologique
franchi toutes les dix minutes n'est plus un niveau psychologique.
\end{insightbox}

Le mécanisme est simple~: le prix de l'or a plus que doublé sur la période et
sa volatilité absolue a suivi, tandis que le pas de la grille --- dix dollars
--- est resté ce qu'il est. L'\code{ATR} de cinq minutes passe de
\XTenAtrFirst{} à \XTenAtrLast{} dollars entre le début et la fin de
l'échantillon.

\subsection{Ce que la dérive du régime change pour le signal}

Trois conséquences, toutes défavorables et toutes vérifiées plus loin dans ce
rapport.

\begin{description}
  \item[La cible change de nature.] La cible d'une position est le barreau
  suivant. Elle valait dix-sept écarts-types de barre en 2019 --- une sortie
  pratiquement hors de portée dans la durée de vie maximale d'une position ---
  et moins de quatre en 2025, c'est-à-dire un déplacement que le bruit
  produit seul. La décomposition par tranche de régime
  (section~\ref{sec:x10_in_sample}) montre effectivement que les deux extrêmes
  sont les pires~: \XTenWorstDollarBandR{}~R quand dix dollars valent plus de
  dix \code{ATR}, et le meilleur résultat --- encore négatif ---
  \XTenBestDollarBandR{}~R dans la tranche intermédiaire.
  \item[Le coût relatif s'effondre, au mauvais moment.] Le spread de
  \XTenSpreadUsd{} représente la moitié d'un \code{ATR} de cinq minutes en
  2019 et un dixième en 2025. Le régime le plus favorable au coût est
  précisément celui où la cible est du bruit~: les deux effets tirent en sens
  contraire et ne se compensent pas.
  \item[Le garde-fou de rapport gain/risque devient mordant.] Le mandat exige
  qu'une position ne soit prise que si le gain visé vaut au moins la perte
  acceptée. Sur un régime où la cible rétrécit en unités de volatilité tandis
  que le stop reste proportionnel à cette volatilité, ce test se met à
  refuser. Il rejetait \XTenRRejectFirstYear{} des candidats en 2019~; il en
  rejette \XTenRRejectLastYear{} en 2025.
\end{description}

% GÉNÉRÉ par scripts/build_x10_report_assets.py — ne pas éditer à la main.
% Toute valeur vient de results/xau_x10/*.json.
\begin{table}[htbp]
\centering
\small
\caption{Taux de rejet du garde-fou $R \geq 1$, par année. Le dénominateur ne compte que les candidats ayant atteint le test~: un refus au contexte n'a pas de $R$ estimé.}
\label{tab:x10_r_rule}
\begin{tabular}{l r r r}
\toprule
\textbf{Année} & \textbf{Candidats au test R} & \textbf{Rejetés} & \textbf{Taux de rejet} \\
\midrule
2019 & $166$ & $11$ & $6{,}6$\,\% \\
2020 & $338$ & $90$ & $26{,}6$\,\% \\
2021 & $274$ & $39$ & $14{,}2$\,\% \\
2022 & $296$ & $49$ & $16{,}6$\,\% \\
2023 & $272$ & $55$ & $20{,}2$\,\% \\
2024 & $412$ & $103$ & $25{,}0$\,\% \\
2025 & $741$ & $412$ & $55{,}6$\,\% \\
\bottomrule
\end{tabular}
\end{table}


\begin{figure}[htbp]
\centering
\includegraphics[width=0.92\linewidth]{x10_fig08_regle_r.png}
\caption{Taux de rejet annuel du garde-fou de rapport gain/risque. Le
dénominateur ne compte que les candidats ayant atteint ce test~: un candidat
refusé plus tôt, au contexte, n'a pas de rapport estimé et n'y figure pas.}
\end{figure}

\begin{warningbox}
\textbf{Ce filtre avait été jugé inopérant avant d'être compté, et ce jugement
était faux.} Le raisonnement initial supposait une entrée \emph{au niveau}~;
l'entrée réelle n'a jamais lieu au niveau. Pour un \emph{breakout}, elle suit
une période de maintien, pendant laquelle le prix a déjà parcouru une part de
la cible. Pour un \emph{reversal}, elle suit une excursion au-delà du niveau,
dont la profondeur éloigne le stop. Dans les deux cas, le rapport effectif est
structurellement inférieur au rapport théorique, et l'écart grandit avec la
volatilité --- d'où le décrochage de la dernière année. L'erreur est conservée
et signalée dans les notes de recherche~: raisonner sur une entrée idéalisée
avant d'avoir un moteur est une méthode qui se paie.
\end{warningbox}

\subsection{Ce que cette dérive interdit de conclure}

Deux interdictions, posées avant la mesure et tenues après elle.

D'abord, on ne peut pas \textbf{remplacer la grille en dollars par une grille
en multiples d'\code{ATR}} au vu de ce tableau. Ce serait une autre stratégie
que celle commandée~: une autre spécification, un autre budget d'essais, et
surtout une décision prise après avoir vu les données, ce qui est précisément
ce que le dispositif de cette phase cherche à empêcher.

Ensuite, on ne peut pas \textbf{lire la moindre tendance temporelle du
résultat comme une amélioration}. La perte s'atténue au fil des années
(section~\ref{sec:x10_in_sample}), et deux explications purement mécaniques
suffisent à en rendre compte --- un spread constant qui pèse de moins en moins
en proportion d'un prix qui monte, et un garde-fou qui écarte de plus en plus
de candidats. Cette campagne ne les sépare pas, et aucune des deux ne
constitue un signal.
